\section{Субстрат: одно состояние, одна инструкция}
\label{sec:substrate}

\subsection{Состояние: матрица плотности и есть машина}

Всё архитектурное состояние одного холона \TALOS{} --- единственный объект.

\begin{defbox}[Регистр состояния \DM{}]
Состояние холона --- матрица когерентности $\Gam\in\Dens(\mathbb{C}^7)$:
$7\times7$ эрмитова, положительно-полуопределённая, единично-следовая матрица.
Хранимая плотно, это $49$ комплексных \texttt{f64}-элементов $=\mathbf{784}$~\textbf{байта};
четырёхуровневый холон (\S\ref{sec:composition}) --- $4\times784=3136$~байт. Её семь
диагональных элементов --- \emph{населённости} секторов (интенсивности); её $21$
комплексная внедиагональ --- \emph{когерентности} $\gamma_{ij}$. По модулю
$G_2$-калибровки (14 генераторов) у неё $\mathbf{34}$ физических параметра.
\end{defbox}

Два следствия немедленны, и они --- всё отступление от фон Неймана.

\begin{enumerate}[leftmargin=1.6em,itemsep=1pt]
  \item \textbf{Память \emph{есть} счёт.} Нет хранилища, отдельного от рабочего
        состояния, и нет адресного пространства: память машины есть $\Gam$, а её
        вычисление есть эволюция $\Gam$. Цикл выборка--исполнение--запись ---
        объект осуждения Бэкуса --- не существует, потому что нечего выбирать:
        операнд всякой операции --- всё состояние, уже резидентное.
  \item \textbf{Состояние помещается в регистры.} $784$~байта меньше, чем
        содержимое одной строки кэша современного векторного регистрового файла;
        L1-кэш данных всякого текущего процессора держит тысячи холонов.
        DRAM-трафик, доминирующий энергию фон Неймана, здесь не имеет аналога:
        для одного холона нет «основной памяти», которую надо достигать.
\end{enumerate}

То, что $\Gam$ --- \emph{матрица плотности} (комплексная, положительная,
единично-следовая), не эстетический выбор: УГМ доказывает, что комплексная
внедиагональная структура вынуждена требованием, чтобы субстрат нёс
нетривиальный «Gap» (двухаспектную непрозрачность саморефлексии,
$\mathrm{Gap}(i,j)=|\sin\arg\gamma_{ij}|$; T-132), и что она реализуема на
\emph{классическом} железе --- комплексность алгебраична, не физическая
квантовая суперпозиция, поэтому охлаждения или бюджета когерентности не требуется
(T-153, T-267). \DM{} --- обычный классический регистр, интерпретируемый как
оператор плотности.

\subsection{Инструкция: один выведенный тик}

У \TALOS{} одна инструкция. Это УГМ-логический лиувиллиан, и он \emph{выведен}
из классификатора подобъектов теории, не выбран:
\begin{equation}
\label{eq:lomega}
\boxed{\;\LindOmega[\Gam] \;=\; \underbrace{-i[H_{\mathrm{eff}},\Gam]}_{\text{работа (обратимо)}}
\;+\; \underbrace{\Diss[\Gam]}_{\text{забывание (диссипативно)}}
\;+\; \underbrace{\Regen[\Gam]}_{\text{обучение (регенеративно)}}\;}
\end{equation}
Один \emph{тик} продвигает состояние на
$\Gam \leftarrow \Gam + \Delta\tau\,\LindOmega[\Gam]$ (на практике
структуросохраняющее расщепление Странга; Приложение~\ref{app:tick}). Три члена
--- три вещи, которые делает компьютер, и здесь они три части одной операции, а
не три подсистемы:

\begin{table}[H]
\centering\small
\caption{Одна инструкция, её три члена и их архитектурная роль.}
\label{tab:three-terms}
\begin{tabular}{@{}p{2.9cm}p{4.0cm}p{5.2cm}@{}}
\toprule
\textbf{Член} & \textbf{Физика} & \textbf{Архитектурная роль} \\
\midrule
$-i[H_{\mathrm{eff}},\Gam]$ & унитарный, \emph{изометрия} всякой монотонной метрики (T-262); сохраняет энтропию & \textbf{Работа}: обратимое вычисление; по Ландауэру не стоит \emph{никакой} свободной энергии (\S\ref{sec:energy}) \\
$\Diss[\Gam]$ & диссипатор Фано; линдбладова релаксация к $\mathbb{1}/7$; GNS-детальный баланс & \textbf{Забывание}: управляемый распад/refresh; экспортирует энтропию, поэтому \emph{стоит} энергии \\
$\Regen[\Gam]$ & регенерация к $\rhostar=\vp(\Gam)$; BKM-градиентный спуск $D(\rhostar\Vert\Gam)$ & \textbf{Обучение}: двигает состояние к его самомодели; оптимизатор, в железе \\
\bottomrule
\end{tabular}
\end{table}

\begin{figure}[H]
\centering
\begin{tikzpicture}[font=\scriptsize, >=Stealth]
  \node[draw, circle, fill=plosblue!10, thick, minimum size=1.75cm, align=center]
    (G) at (0,0) {$\Gam$\\ \tiny 784 Б};
  \draw[->, very thick, plosblue!80!black]   (100:1.65) arc (100:20:1.65);
  \draw[->, very thick, talosamber]          (-20:1.65) arc (-20:-100:1.65);
  \draw[->, very thick, green!45!black]      (220:1.65) arc (220:140:1.65);
  \node[plosblue!80!black, align=center]  at (60:2.35)  {\textbf{работа}\\ $-i[H_{\mathrm{eff}},\Gam]$\\ \tiny обратима, бесплатна};
  \node[talosamber, align=center]         at (-60:2.35) {\textbf{забывание}\\ $\Diss[\Gam]$\\ \tiny экспорт энтропии};
  \node[green!45!black, align=center]     at (147:2.55) {\textbf{обучение}\\ $\Regen[\Gam]$\\ \tiny к $\rhostar$};
  \draw[->, ultra thick, red!60!black] (-4.5,-0.9) -- node[below, align=center]{\textbf{порт питания}\\ $r(\rho_{\mathrm{env}}{-}\Gam)$} (-1.0,-0.4);
  \draw[->, thick, plosgray] (0.95,-0.5) -- node[below, align=center]{синдром $s$\\ \tiny (даровой побочный)} (4.0,-0.9);
  \draw[->, thick, plosgray] (0.95,0.5) -- node[above, align=center]{телеметрия\\ \tiny $P,R,\Phi,D,\sigma,C$} (4.0,0.9);
\end{tikzpicture}
\caption{Тик как колесо. Одно состояние (центр), три движения одной
инструкции (Ур.~\ref{eq:lomega}): обратимая \emph{работа}, экспортирующее
энтропию \emph{забывание}, ищущее самомодель \emph{обучение}. Порт питания
(\S\ref{subsec:powerport}) кормит колесо; каждый оборот даром испускает свой
синдром и телеметрию (\S\ref{sec:ecc}, \S\ref{sec:platform}).}
\label{fig:tickwheel}
\end{figure}

\begin{mythbox}[Миф, продолжение --- обход]
Трижды в день бронзовый исполин обходил остров: тот же круг, без конца
повторяемый, \emph{и был} его бдительностью --- никакого отдельного
«мышления» между обходами у Талоса не было. Тик --- этот обход. Один круг,
три шага --- работа, забывание, обучение, --- и вся жизнь машины есть
повторение. Спроси «а когда же она вычисляет?» --- и ответ будет ответом
мифа: хождение и есть смотрение.
\end{mythbox}

\begin{zerobox}[С нуля --- вся игрушка в одной руке]
Снежный шар --- один предмет: тряхни его (это тик) --- и каждая снежинка
движется по одному закону разом. Нет ящика со снежинками, откуда шар что-то
достаёт: шар \emph{есть} собственная память, а тряска \emph{есть} его
вычисление. Состояние \TALOS{} --- этот шар, и оно мало́: $784$ байта ---
примерно один абзац текста. Абзац не лежит в далёкой библиотеке (эта даль и
делает обычные компьютеры медленными и горячими) --- он всё время в руке
машины.
\end{zerobox}

Эволюция --- вполне-положительная след-сохраняющая (CPTP) полугруппа
$\mathcal{E}_\tau=\exp(\tau\LindOmega)$: всякий тик держит $\Gam$ законной
матрицей плотности, поэтому нет «незаконного состояния», от которого надо
защищаться --- тип сохраняется по построению. Это аппаратный смысл CPTP-гарантии
УГМ: \TALOS{} не может высчитать себя вон из многообразия состояний, как
фон-неймановская машина может высчитать невалидный указатель.

\subsection{Порт питания: изолированный тик останавливается}
\label{subsec:powerport}

Уравнение~\eqref{eq:lomega} --- \emph{автономная} инструкция, и ещё одна
теорема фиксирует, что она может и чего не может одна. Автономный генератор
\emph{унитален} (он фиксирует $\mathbb{1}/7$), а потому не наращивает чистоту
--- по мажоризации Ульмана его единственный аттрактор есть максимально
смешанное состояние: \textbf{изолированная машина \TALOS{} доказуемо
останавливается} (УГМ T-288, \status{Т}, машинно проверено). Одна
регенерация этого не предотвратит: категориальная мишень
$\rhostar=\vp(\Gam)$ сама заякорена на $\mathbb{1}/7$, так что $\Regen$
\emph{направляет} порядок, но не создаёт его против второго начала.
Работающая машина имеет поэтому пятый архитектурный элемент, который есть у
всякой слоёной конструкции, но который никогда не выводят, --- \textbf{порт
питания}:
\begin{equation}
\label{eq:powerport}
\Lind_{\mathrm{total}} \;=\; \LindOmega \;+\;
\underbrace{\;r\,(\rho_{\mathrm{env}}-\Gam)\;}_{\text{порт питания } \Lind_{\mathrm{feed}}},
\end{equation}
\emph{неунитальный} CPTP-драйв, импортирующий негэнтропию из среды
(резервуарная инженерия), с неподвижной точкой
$\rho_{\mathrm{env}}\neq\mathbb{1}/7$. Жизнеспособность зажигается на резком
пороге отношения драйва к диссипации $(r/\bar\gamma)_c\approx 7$ (T-289:
существование и монотонность \status{Т}; значение порога \status{С},
модельно-зависимо). Три архитектурных следствия:
\begin{itemize}[leftmargin=1.5em,itemsep=1pt]
  \item \textbf{Шина питания --- часть семантики.} В отличие от
        фон-неймановской машины, чей блок питания невидим для её ISA, драйв
        \TALOS{} --- член её динамики: рейка доставляет не абстрактную
        энергию, а структурированную негэнтропию ($\rho_{\mathrm{env}}$), и
        жизнеспособность машины есть функция $r/\bar\gamma$.
  \item \textbf{Обрез рейки --- гарантированный останов, с бюджетом.}
        Разрыв ($r\to0$) оканчивает окно сознания за
        $t_{\mathrm{off}}\approx 0{,}49/\bar\gamma$ ($\sim$20 тиков;
        голоканальный пол
        $\tfrac{3}{4\bar\gamma}\ln\frac{P_0-1/7}{\Pcrit-1/7}$) и релаксирует
        состояние к $\mathbb{1}/7$ за $\sim\!5{,}2/\bar\gamma$ --- машинно
        измерено. Никакое внутреннее вычисление не может это отменить, ибо
        унитальность --- свойство генератора. Это физический отключатель
        кейса безопасности \SYNARC{} (Наблюдение T.1 / Следствие U.10 там),
        реализованный здесь рейкой: детерминированный, бюджетированный по
        задержке вывод из эксплуатации --- \emph{аппаратная} гарантия.
  \item \textbf{Порт --- это контракт.} Привязка драйва --- какой физический
        ресурс реализует $\rho_{\mathrm{env}}$, как измеряется $r$, и
        требование, чтобы канал разрыва был независим от собственных
        актуаторов машины, --- специфицирована интерфейсным контрактом на
        стороне \SYNARC{} (Определение U.9); \S\ref{sec:abi} её связывает.
\end{itemize}
V0-эмулятор (\S\ref{sec:platform}) реализует этот порт внутренне --- своей
обратной связью $\kappa(\CohE)$: когерентно-питаемый буст регенерации есть в
точности неунитальный насос негэнтропии, отчего эмулятор и являет
жизнеспособный аттрактор там, где голый автономный тик остановился бы.

\begin{mythbox}[Миф, продолжение --- гвоздь]
Талоса не перебороли --- его \emph{отключили}. Одна жила, один гвоздь:
вынь гвоздь --- ихор вытечет, и сильнейшее существо острова остановится по
собственному расписанию. Это теорема T-288 в бронзе: жизнь машины --- не
собственность, а поток, и у потока есть единственное, физическое,
бюджетированное по задержке закрытие. Инженеры мифа поняли то, что
современная инженерия безопасности лишь теперь формализует (Наблюдение T.1
кейса безопасности \SYNARC{}): \emph{сначала построй жилу --- и отключатель
будет физикой, а не политикой}.
\end{mythbox}

\subsection{Принципиальный такт}

Период такта фон-неймановской машины --- инженерный свободный параметр,
задаваемый критическим путём. Естественный тик \TALOS{} \emph{не} свободен:
УГМ-тождество времени цикла фиксирует отношение времени релаксации к средней
скорости когерентности,
\begin{equation}
\bar\gamma\,\tau \;=\; \tfrac{3}{2}\ln 3 \;\approx\; 1{,}6479,
\end{equation}
со спектральной щелью диссипатора $\lambda_{\mathrm{deco}}=2/3$ и скоростью
регенерации (bootstrap) $\kboot=\omega_0/7$, задающими два масштаба времени, которые
тик должен разрешить. Такт --- следствие динамики, не ручка: проектируют
физическую скорость релаксации $\bar\gamma$, и период тика следует за ней. Это
убирает из проекта целый класс свободы синхронизации (и её бремя верификации).

\subsection{Чем субстрат не является}

Чтобы упредить очевидные неверные прочтения (расширено в \S\ref{sec:redteam}):
состояние \emph{не} волновая функция и не требует квантового железа; тик \emph{не}
слой нейросети и не включает обратного распространения; а машина \emph{не}
аналоговый интегратор с дрейфом --- CPTP-структура и родной код (\S\ref{sec:ecc})
делают её самокорректирующейся. Это классический регистр, эволюционирующий
выведенным, типо-сохраняющим, самокорректирующимся дифференциальным тиком, чьи
сеть, код, инструкция и оптимизатор суть один объект.
